\section{Related Work}
\label{sec:related}

% MANUSCRIPT-NODE: RELATED-P01
Quality-aware policy fitting already provides several ways to use or suppress suboptimal behavior. Advantage-weighted regression, implicit value learning, critic-regularized regression, and proximal behavior objectives emphasize actions estimated to be better while reducing the impact of lower-quality data~\cite{peng2019advantage,kostrikov2021offline,wang2020critic,zhuang2023behavior}. Positive-only and hard-filtering variants take the conservative endpoint and remove selected negative updates altogether. At the same time, failed or suboptimal behavior can remain informative: it can suppress competing bad modes, sharpen a decision boundary, or release probability mass for alternatives represented by the model. The established picture is therefore not that negative feedback is simply noise, but that it combines information value with optimization risk. What remains under-characterized is the temporal relation to the current learner. A fixed negative action can begin as relevant local feedback and become increasingly remote after repeated reuse. We study that transition and its aggregate consequences, complementing rather than dismissing prior weighting and filtering methods.
% END-MANUSCRIPT-NODE: RELATED-P01

% MANUSCRIPT-NODE: RELATED-P02
Off-policy learning has long controlled mismatch between the behavior distribution and the current policy. Importance correction, PPO-style clipping, and trust regions limit the effect of samples collected by another policy~\cite{schulman2017proximal}. Replay-based algorithms and maximum-entropy off-policy control likewise manage repeated data reuse while the learner evolves~\cite{haarnoja2018soft}. Offline-RL analyses make the same mismatch explicit when no new interaction is available~\cite{levine2020offline,fu2020d4rl}. Low action probability or high surprisal can therefore serve as a useful learner-relative warning signal. These approaches establish that mismatch and rarity matter; our distinction concerns how they arise. We treat remoteness not only as a static property used to choose a weight, but as an endogenous state variable that a negative update increases and subsequent reuse feeds back into the next update. We then connect that per-sample process to aggregate equilibria and targeted causal interventions. This perspective complements clipping and stale-policy control rather than asserting that they are ineffective.
% END-MANUSCRIPT-NODE: RELATED-P02

% MANUSCRIPT-NODE: RELATED-P03
Offline reinforcement learning contains several distinct failure channels, and far-field signed actor dynamics is only one of them. Conservative Q-learning controls value extrapolation by assigning pessimistic values to unsupported actions~\cite{kumar2020conservative}. Implicit Q-learning avoids querying out-of-distribution actions during value improvement and then fits an actor to supported high-value behavior~\cite{kostrikov2021offline}. Behavior-constrained and proximal approaches instead keep the learned policy close to the data or a behavior reference~\cite{fujimoto2019off,wang2020critic,zhuang2023behavior}. Data selection adds another axis by changing which transitions participate in learning. DRPO focuses on the signed actor field after an advantage-like signal has been supplied: it asks how negative actor updates change as historical actions become remote, and how to preserve their local information without allowing their far-field contribution to dominate. This actor-side control can be combined with pessimistic critics, in-support value learning, or behavior regularization. It is not presented as a replacement for those safeguards or as a complete explanation of every offline-RL failure.
% END-MANUSCRIPT-NODE: RELATED-P03
